\section{Conclusion and Future Work}
\label{sec:conclusion}

Obstacles and narrow connections in complex construction maps prevent a swing-arm troweling robot from obtaining a continuous, executable coverage route from either a fixed tool footprint or a static reachability region alone. This paper presented a hierarchical CPP method for known two-dimensional environments with obstacles. Bidirectional rectangular candidates, hole filling, and connectivity completion represent the workable area; a feasible graph organizes global visitation; and first visits and repeated transit visits are instantiated as coordinated base--arm work motions and transfer motions. Planning and evaluation preserve the geometry of the base, arm, and circular disc, and coverage and nominal execution time are calculated from the disc's actual sweep along the continuous trajectory. The method thus connects map decomposition, region visitation, and end-tool coverage verification in one complete pipeline.

On \texttt{simple} and \texttt{complex}, the proposed method achieves full-trajectory coverage ratios of 94.11\% and 95.73\%, respectively, and both routes pass the path check. Although the fixed-width Fields2Cover baseline generates shorter paths, its connecting segments leave the valid region by 4.51~m$^2$ and 20.00~m$^2$ on the two maps, so the associated nominal times cannot be treated as executable-route results. The comparison demonstrates the practical value of adjustable relative motion between base travel and disc operation in dense obstacle fields. In the five-map scan-strategy study, row-then-column candidates give the highest mean full-trajectory coverage (94.32\%) and the lowest mean execution time (2316.60~s), indicating that complementary row and column directions reduce representation bias; a single-direction strategy can nevertheless be preferable in some regular scenes. The speed study further shows that increasing base speed relative to arm angular velocity enlarges the spacing between sweeps and creates a coverage--time trade-off. Route geometry and execution parameters must therefore be selected together.

The current method relies on heuristic decomposition and visitation ordering and does not guarantee global optimality. It assumes a known static two-dimensional map, and its evaluation does not yet include complete chassis dynamics, localization error, disc pressure, or surface-quality feedback. Future work will jointly optimize decomposition and visitation, introduce online perception and local replanning, and validate trajectory execution and finishing quality on a physical swing-arm troweling robot. The framework may also serve as a decision-support tool for construction-map partitioning and base--arm motion configuration.
